% =============================================================
%  04_soluzioni.tex — Soluzioni (148 esercizi)
% =============================================================
\chapter*{Soluzioni}
\addcontentsline{toc}{chapter}{Soluzioni}
\markboth{Soluzioni}{Soluzioni}
% Sezioni non numerate per evitare conflitti con Esercizi nell'indice

\noindent\small\color{black!65}
\unaStella/\dueStelle: solo codice.\enspace
\treStelle: commenti inline.\enspace
\quattroStelle/\cinqueStelle: box ragionamento + codice commentato.
\color{black}

\vspace{8pt}

% =============================================================
\section*{Fondamenti}
% =============================================================

\label{sol:1}
\begin{solbox}{1}
\small\color{black!55}\hyperref[ex:1]{$\leftarrow$ Esercizio}\normalsize\color{black}
\begin{lstlisting}
(defun add (a b) (+ a b))
\end{lstlisting}
\end{solbox}

\label{sol:2}
\begin{solbox}{2}
\small\color{black!55}\hyperref[ex:2]{$\leftarrow$ Esercizio}\normalsize\color{black}
\begin{lstlisting}
(defun circle-area (r)
  (* pi r r))

(defun circle-circumference (r)
  (* 2 pi r))
\end{lstlisting}
\end{solbox}

\label{sol:3}
\begin{solbox}{3}
\small\color{black!55}\hyperref[ex:3]{$\leftarrow$ Esercizio}\normalsize\color{black}
\begin{lstlisting}
(defun celsius->fahrenheit (c)
  (+ (* c 9/5) 32))

(defun fahrenheit->celsius (f)
  (* (- f 32) 5/9))
\end{lstlisting}
\end{solbox}

\label{sol:4}
\begin{solbox}{4}
\small\color{black!55}\hyperref[ex:4]{$\leftarrow$ Esercizio}\normalsize\color{black}
\begin{lstlisting}
(defun ipotenusa (a b)
  (let ((a2 (* a a))
        (b2 (* b b)))
    (sqrt (+ a2 b2))))
\end{lstlisting}
\end{solbox}

\label{sol:5}
\begin{solbox}{5}
\small\color{black!55}\hyperref[ex:5]{$\leftarrow$ Esercizio}\normalsize\color{black}
\begin{lstlisting}
(defun distanza-euclidea (x1 y1 x2 y2)
  (let ((dx (- x2 x1))
        (dy (- y2 y1)))
    (sqrt (+ (* dx dx) (* dy dy)))))
\end{lstlisting}
\end{solbox}

\label{sol:6}
\begin{solbox}{6}
\small\color{black!55}\hyperref[ex:6]{$\leftarrow$ Esercizio}\normalsize\color{black}
\begin{lstlisting}
(defun pari-p (n)   (= (mod n 2) 0))
(defun dispari-p (n) (= (mod n 2) 1))
\end{lstlisting}
\end{solbox}

\label{sol:7}
\begin{solbox}{7}
\small\color{black!55}\hyperref[ex:7]{$\leftarrow$ Esercizio}\normalsize\color{black}
\begin{lstlisting}
(defun divisibile-p (a b)
  (= (mod a b) 0))
\end{lstlisting}
\end{solbox}

\label{sol:8}
\begin{solbox}{8}
\small\color{black!55}\hyperref[ex:8]{$\leftarrow$ Esercizio}\normalsize\color{black}
\begin{lstlisting}
(defun segno (n)
  (cond ((< n 0) 'negativo)
        ((= n 0) 'zero)
        (t       'positivo)))
\end{lstlisting}
\end{solbox}

\label{sol:9}
\begin{solbox}{9}
\small\color{black!55}\hyperref[ex:9]{$\leftarrow$ Esercizio}\normalsize\color{black}
\begin{lstlisting}
(defun max-tre (a b c)
  (cond ((and (>= a b) (>= a c)) a)
        ((>= b c)                 b)
        (t                        c)))
\end{lstlisting}
\end{solbox}

\label{sol:10}
\begin{solbox}{10}
\small\color{black!55}\hyperref[ex:10]{$\leftarrow$ Esercizio}\normalsize\color{black}
\begin{lstlisting}
(defun triangolo-valido-p (a b c)
  (and (< a (+ b c))
       (< b (+ a c))
       (< c (+ a b))))
\end{lstlisting}
\end{solbox}

\label{sol:11}
\begin{solbox}{11}
\small\color{black!55}\hyperref[ex:11]{$\leftarrow$ Esercizio}\normalsize\color{black}
\begin{lstlisting}
(defun bisestile-p (anno)
  (or (and (= (mod anno 4)   0)
           (/= (mod anno 100) 0))
      (= (mod anno 400) 0)))
\end{lstlisting}
\end{solbox}

\label{sol:12}
\begin{solbox}{12}
\small\color{black!55}\hyperref[ex:12]{$\leftarrow$ Esercizio}\normalsize\color{black}
\begin{lstlisting}
(defun giorni-mese (mese anno)
  (cond ((member mese '(1 3 5 7 8 10 12)) 31)
        ((member mese '(4 6 9 11))        30)
        ((bisestile-p anno)               29)
        (t                                28)))
\end{lstlisting}
\end{solbox}

\label{sol:13}
\begin{solbox}{13}
\small\color{black!55}\hyperref[ex:13]{$\leftarrow$ Esercizio}\normalsize\color{black}
\begin{lstlisting}
(defun voto->giudizio (p)
  (cond ((>= p 90) "ottimo")
        ((>= p 75) "buono")
        ((>= p 60) "sufficiente")
        (t         "insufficiente")))
\end{lstlisting}
\end{solbox}

\label{sol:14}
\begin{solbox}{14}
\small\color{black!55}\hyperref[ex:14]{$\leftarrow$ Esercizio}\normalsize\color{black}
\begin{lstlisting}
(defun terna-pitagorica-p (a b c)
  (= (* c c) (+ (* a a) (* b b))))
\end{lstlisting}
\end{solbox}

\label{sol:15}
\begin{solbox}{15}
\small\color{black!55}\hyperref[ex:15]{$\leftarrow$ Esercizio}\normalsize\color{black}
\begin{lstlisting}
(defun radici-quadratiche (a b c)
  (let* ((delta (- (* b b) (* 4 a c))))
    (if (< delta 0)
        nil
        (let* ((sd (sqrt delta))
               (r1 (/ (+ (- b) sd) (* 2 a)))
               (r2 (/ (- (- b) sd) (* 2 a))))
          (list r1 r2)))))
\end{lstlisting}
\end{solbox}

\label{sol:16}
\begin{solbox}{16}
\small\color{black!55}\hyperref[ex:16]{$\leftarrow$ Esercizio}\normalsize\color{black}
\begin{lstlisting}
(defun calcola (op a b)
  (case op
    (+  (+ a b))
    (-  (- a b))
    (*  (* a b))
    (/  (if (= b 0) nil (/ a b)))))
\end{lstlisting}
\end{solbox}

\label{sol:17}
\begin{solbox}{17}
\small\color{black!55}\hyperref[ex:17]{$\leftarrow$ Esercizio}\enspace\textcolor{black!30}{|}\enspace\hyperref[hint:17]{$\leftarrow$ Suggerimento}\normalsize\color{black}
\begin{lstlisting}
(defun somma-cifre (n)
  (if (= n 0)
      0
      (+ (mod n 10)
         (somma-cifre (floor n 10)))))
\end{lstlisting}
\end{solbox}

\label{sol:18}
\begin{solbox}{18}
\small\color{black!55}\hyperref[ex:18]{$\leftarrow$ Esercizio}\normalsize\color{black}
\begin{lstlisting}
(defun conta-fino-a (n)
  (dotimes (i n)
    (print (1+ i))))
\end{lstlisting}
\end{solbox}

\label{sol:19}
\begin{solbox}{19}
\small\color{black!55}\hyperref[ex:19]{$\leftarrow$ Esercizio}\normalsize\color{black}
\begin{lstlisting}
(defun tabellina (n)
  (dotimes (i 10)
    (format t "~D x ~D = ~D~%"
            n (1+ i) (* n (1+ i)))))
\end{lstlisting}
\end{solbox}

\label{sol:20}
\begin{solbox}{20}
\small\color{black!55}\hyperref[ex:20]{$\leftarrow$ Esercizio}\normalsize\color{black}
\begin{lstlisting}
(defun sconto (prezzo)
  (cond ((>= prezzo 200) (* prezzo 0.8))
        ((>= prezzo 100) (* prezzo 0.9))
        (t               prezzo)))
\end{lstlisting}
\end{solbox}

\label{sol:21}
\begin{solbox}{21}
\small\color{black!55}\hyperref[ex:21]{$\leftarrow$ Esercizio}\normalsize\color{black}
\begin{lstlisting}
(defun paga-settimanale (ore tariffa)
  (if (<= ore 40)
      (* ore tariffa)
      (+ (* 40 tariffa)
         (* (- ore 40) tariffa 3/2))))
\end{lstlisting}
\end{solbox}

% =============================================================
\section*{Liste}
% =============================================================

\label{sol:22}
\begin{solbox}{22}
\small\color{black!55}\hyperref[ex:22]{$\leftarrow$ Esercizio}\normalsize\color{black}
\begin{lstlisting}
;; Parte A: con cons
(cons 'a (cons (cons 'b (cons 'c nil)) (cons 'd nil)))
; -> (A (B C) D)

;; Parte B: con list
(list 'a (list 'b 'c) 'd)
; -> (A (B C) D)
\end{lstlisting}
\end{solbox}

\label{sol:23}
\begin{solbox}{23}
\small\color{black!55}\hyperref[ex:23]{$\leftarrow$ Esercizio}\normalsize\color{black}
\begin{lstlisting}
(car (cdr (cdr '(a b c d e))))
; -> C
\end{lstlisting}
\end{solbox}

\label{sol:24}
\begin{solbox}{24}
\small\color{black!55}\hyperref[ex:24]{$\leftarrow$ Esercizio}\normalsize\color{black}
\begin{lstlisting}
(defun first-element (lst)
  (if (null lst) nil (nth 0 lst)))

(defun rest-elements (lst)
  (if (null lst) nil (nthcdr 1 lst)))
\end{lstlisting}
\end{solbox}

\label{sol:25}
\begin{solbox}{25}
\small\color{black!55}\hyperref[ex:25]{$\leftarrow$ Esercizio}\normalsize\color{black}
\begin{lstlisting}
(defun primo-e-ultimo (lst)
  (if (null lst)
      nil
      (list (car lst) (car (last lst)))))
\end{lstlisting}
\end{solbox}

\label{sol:26}
\begin{solbox}{26}
\small\color{black!55}\hyperref[ex:26]{$\leftarrow$ Esercizio}\normalsize\color{black}
\begin{lstlisting}
(defun scambia-primi-due (lst)
  (if (or (null lst) (null (cdr lst)))
      lst
      (cons (cadr lst)
            (cons (car lst)
                  (cddr lst)))))
\end{lstlisting}
\end{solbox}

\label{sol:27}
\begin{solbox}{27}
\small\color{black!55}\hyperref[ex:27]{$\leftarrow$ Esercizio}\normalsize\color{black}
\begin{lstlisting}
(defun length* (lst)
  (if (null lst)
      0
      (1+ (length* (cdr lst)))))
\end{lstlisting}
\end{solbox}

\label{sol:28}
\begin{solbox}{28}
\small\color{black!55}\hyperref[ex:28]{$\leftarrow$ Esercizio}\normalsize\color{black}
\begin{lstlisting}
(defun sum-list (lst)
  (if (null lst)
      0
      (+ (car lst) (sum-list (cdr lst)))))
\end{lstlisting}
\end{solbox}

\label{sol:29}
\begin{solbox}{29}
\small\color{black!55}\hyperref[ex:29]{$\leftarrow$ Esercizio}\normalsize\color{black}
\begin{lstlisting}
(defun product-list (lst)
  (if (null lst)
      1
      (* (car lst) (product-list (cdr lst)))))
\end{lstlisting}
\end{solbox}

\label{sol:30}
\begin{solbox}{30}
\small\color{black!55}\hyperref[ex:30]{$\leftarrow$ Esercizio}\normalsize\color{black}
\begin{lstlisting}
(defun max-list (lst)
  (cond ((null lst) (error "lista vuota"))
        ((null (cdr lst)) (car lst))
        (t (max (car lst) (max-list (cdr lst))))))
\end{lstlisting}
\end{solbox}

\label{sol:31}
\begin{solbox}{31}
\small\color{black!55}\hyperref[ex:31]{$\leftarrow$ Esercizio}\normalsize\color{black}
\begin{lstlisting}
(defun member* (item lst)
  (cond ((null lst)          nil)
        ((equal (car lst) item) t)
        (t (member* item (cdr lst)))))
\end{lstlisting}
\end{solbox}

\label{sol:32}
\begin{solbox}{32}
\small\color{black!55}\hyperref[ex:32]{$\leftarrow$ Esercizio}\normalsize\color{black}
\begin{lstlisting}
(defun count-element (item lst)
  (cond ((null lst) 0)
        ((equal (car lst) item)
         (1+ (count-element item (cdr lst))))
        (t (count-element item (cdr lst)))))
\end{lstlisting}
\end{solbox}

\label{sol:33}
\begin{solbox}{33}
\small\color{black!55}\hyperref[ex:33]{$\leftarrow$ Esercizio}\normalsize\color{black}
\begin{lstlisting}
(defun position* (item lst)
  (labels ((iter (lst i)
             (cond ((null lst) nil)
                   ((equal (car lst) item) i)
                   (t (iter (cdr lst) (1+ i))))))
    (iter lst 0)))
\end{lstlisting}
\end{solbox}

\label{sol:34}
\begin{solbox}{34}
\small\color{black!55}\hyperref[ex:34]{$\leftarrow$ Esercizio}\normalsize\color{black}
\begin{lstlisting}
(defun last-element (lst)
  (cond ((null lst) nil)
        ((null (cdr lst)) (car lst))
        (t (last-element (cdr lst)))))
\end{lstlisting}
\end{solbox}

\label{sol:35}
\begin{solbox}{35}
\small\color{black!55}\hyperref[ex:35]{$\leftarrow$ Esercizio}\normalsize\color{black}
\begin{lstlisting}
(defun incrementa-tutti (lst)
  (if (null lst)
      nil
      (cons (1+ (car lst))
            (incrementa-tutti (cdr lst)))))
\end{lstlisting}
\end{solbox}

\label{sol:36}
\begin{solbox}{36}
\small\color{black!55}\hyperref[ex:36]{$\leftarrow$ Esercizio}\normalsize\color{black}
\begin{lstlisting}
(defun solo-pari (lst)
  (cond ((null lst) nil)
        ((evenp (car lst))
         (cons (car lst) (solo-pari (cdr lst))))
        (t (solo-pari (cdr lst)))))
\end{lstlisting}
\end{solbox}

\label{sol:37}
\begin{solbox}{37}
\small\color{black!55}\hyperref[ex:37]{$\leftarrow$ Esercizio}\normalsize\color{black}
\begin{lstlisting}
(defun elimina-tutte (item lst)
  (cond ((null lst) nil)
        ((equal (car lst) item)
         (elimina-tutte item (cdr lst)))
        (t (cons (car lst)
                 (elimina-tutte item (cdr lst))))))
\end{lstlisting}
\end{solbox}

\label{sol:38}
\begin{solbox}{38}
\small\color{black!55}\hyperref[ex:38]{$\leftarrow$ Esercizio}\normalsize\color{black}
\begin{lstlisting}
(defun seleziona-dispari (lst)
  (cond ((null lst)       nil)
        ((null (cdr lst)) (list (car lst)))
        (t (cons (car lst)
                 (seleziona-dispari (cddr lst))))))
\end{lstlisting}
\end{solbox}

\label{sol:39}
\begin{solbox}{39}
\small\color{black!55}\hyperref[ex:39]{$\leftarrow$ Esercizio}\normalsize\color{black}
\begin{lstlisting}
;; Versione ricorsiva semplice (non tail-recursive)
(defun reverse*-naive (lst)
  (if (null lst)
      nil
      (append (reverse*-naive (cdr lst))
              (list (car lst)))))

;; Versione con accumulatore (tail-recursive)
(defun reverse* (lst)
  (labels ((iter (lst acc)
             (if (null lst)
                 acc
                 (iter (cdr lst) (cons (car lst) acc)))))
    (iter lst nil)))
\end{lstlisting}
\end{solbox}

\label{sol:40}
\begin{solbox}{40}
\small\color{black!55}\hyperref[ex:40]{$\leftarrow$ Esercizio}\enspace\textcolor{black!30}{|}\enspace\hyperref[hint:40]{$\leftarrow$ Suggerimento}\normalsize\color{black}
\begin{lstlisting}
(defun append* (lst1 lst2)
  (if (null lst1)
      lst2
      (cons (car lst1)
            (append* (cdr lst1) lst2))))
\end{lstlisting}
\end{solbox}

\label{sol:41}
\begin{solbox}{41}
\small\color{black!55}\hyperref[ex:41]{$\leftarrow$ Esercizio}\enspace\textcolor{black!30}{|}\enspace\hyperref[hint:41]{$\leftarrow$ Suggerimento}\normalsize\color{black}
\begin{lstlisting}
(defun remove-first (item lst)
  (cond ((null lst) nil)
        ((equal (car lst) item) (cdr lst))
        (t (cons (car lst)
                 (remove-first item (cdr lst))))))
\end{lstlisting}
\end{solbox}

\label{sol:42}
\begin{solbox}{42}
\small\color{black!55}\hyperref[ex:42]{$\leftarrow$ Esercizio}\normalsize\color{black}
\begin{lstlisting}
(defun replace-first (old new lst)
  (cond ((null lst) nil)
        ((equal (car lst) old)
         (cons new (cdr lst)))
        (t (cons (car lst)
                 (replace-first old new (cdr lst))))))
\end{lstlisting}
\end{solbox}

\label{sol:43}
\begin{solbox}{43}
\small\color{black!55}\hyperref[ex:43]{$\leftarrow$ Esercizio}\normalsize\color{black}
\begin{lstlisting}
(defun palindroma-p (lst)
  (equal lst (reverse* lst)))
\end{lstlisting}
\end{solbox}

\label{sol:44}
\begin{solbox}{44}
\small\color{black!55}\hyperref[ex:44]{$\leftarrow$ Esercizio}\normalsize\color{black}
\begin{lstlisting}
(defun range (a b)
  (if (> a b)
      nil
      (cons a (range (1+ a) b))))
\end{lstlisting}
\end{solbox}

\label{sol:45}
\begin{solbox}{45}
\small\color{black!55}\hyperref[ex:45]{$\leftarrow$ Esercizio}\normalsize\color{black}
\begin{lstlisting}
(defun duplica-elementi (lst)
  (if (null lst)
      nil
      (cons (car lst)
            (cons (car lst)
                  (duplica-elementi (cdr lst))))))
\end{lstlisting}
\end{solbox}

\label{sol:46}
\begin{solbox}{46}
\small\color{black!55}\hyperref[ex:46]{$\leftarrow$ Esercizio}\normalsize\color{black}
\begin{lstlisting}
(defun schiaccia-coppia (alist)
  (mapcar #'car alist))
\end{lstlisting}
\end{solbox}

\label{sol:47}
\begin{solbox}{47}
\small\color{black!55}\hyperref[ex:47]{$\leftarrow$ Esercizio}\normalsize\color{black}
\begin{lstlisting}
(defun prefisso-p (sub lst)
  (cond ((null sub) t)
        ((null lst)  nil)
        ((equal (car sub) (car lst))
         (prefisso-p (cdr sub) (cdr lst)))
        (t nil)))
\end{lstlisting}
\end{solbox}

\label{sol:48}
\begin{solbox}{48}
\small\color{black!55}\hyperref[ex:48]{$\leftarrow$ Esercizio}\normalsize\color{black}
\begin{lstlisting}
(defun intercala (lst1 lst2)
  (if (null lst1)
      nil
      (cons (car lst1)
            (cons (car lst2)
                  (intercala (cdr lst1) (cdr lst2))))))
\end{lstlisting}
\end{solbox}

\label{sol:49}
\begin{solbox}{49}
\small\color{black!55}\hyperref[ex:49]{$\leftarrow$ Esercizio}\enspace\textcolor{black!30}{|}\enspace\hyperref[hint:49]{$\leftarrow$ Suggerimento}\normalsize\color{black}
\begin{lstlisting}
(defun rle-encode (lst)
  (if (null lst)
      nil
      (labels ((iter (lst curr count acc)
                 (cond
                   ((null lst)
                    (reverse (cons (list count curr) acc)))
                   ((equal (car lst) curr)
                    (iter (cdr lst) curr (1+ count) acc))
                   (t
                    (iter (cdr lst) (car lst) 1
                          (cons (list count curr) acc))))))
        (iter (cdr lst) (car lst) 1 nil))))

(defun rle-decode (lst)
  (if (null lst)
      nil
      (let ((n (caar lst))
            (e (cadar lst)))
        (append (loop repeat n collect e)
                (rle-decode (cdr lst))))))
\end{lstlisting}
\end{solbox}

\label{sol:50}
\begin{solbox}{50}
\small\color{black!55}\hyperref[ex:50]{$\leftarrow$ Esercizio}\enspace\textcolor{black!30}{|}\enspace\hyperref[hint:50]{$\leftarrow$ Suggerimento}\normalsize\color{black}
\begin{lstlisting}
(defun remove-duplicates* (lst)
  (labels ((iter (lst seen)
             (cond
               ((null lst) nil)
               ((member (car lst) seen)
                (iter (cdr lst) seen))
               (t (cons (car lst)
                        (iter (cdr lst)
                              (cons (car lst) seen)))))))
    (iter lst nil)))
\end{lstlisting}
\end{solbox}

% =============================================================
\section*{Ricorsione avanzata}
% =============================================================

\label{sol:51}
\begin{solbox}{51}
\small\color{black!55}\hyperref[ex:51]{$\leftarrow$ Esercizio}\normalsize\color{black}
\begin{lstlisting}
(defun fattoriale (n)
  ;; labels definisce iter localmente (tail-recursive)
  (labels ((iter (n acc)
             (if (= n 0)
                 acc                        ; caso base: restituisci acc
                 (iter (1- n) (* n acc))))) ; accumula n nel prodotto
    (iter n 1)))                            ; avvia con acc = 1
\end{lstlisting}
\end{solbox}

\label{sol:52}
\begin{solbox}{52}
\small\color{black!55}\hyperref[ex:52]{$\leftarrow$ Esercizio}\normalsize\color{black}
\begin{lstlisting}
(defun reverse-tr (lst)
  ;; l'accumulatore accumula in ordine inverso
  (labels ((iter (lst acc)
             (if (null lst)
                 acc
                 (iter (cdr lst) (cons (car lst) acc)))))
    (iter lst nil)))
\end{lstlisting}
\end{solbox}

\label{sol:53}
\begin{solbox}{53}
\small\color{black!55}\hyperref[ex:53]{$\leftarrow$ Esercizio}\normalsize\color{black}
\begin{lstlisting}
(defun somma-range (a b)
  (labels ((iter (i acc)
             (if (> i b)
                 acc
                 (iter (1+ i) (+ acc i)))))
    (iter a 0)))
\end{lstlisting}
\end{solbox}

\label{sol:54}
\begin{solbox}{54}
\small\color{black!55}\hyperref[ex:54]{$\leftarrow$ Esercizio}\enspace\textcolor{black!30}{|}\enspace\hyperref[hint:54]{$\leftarrow$ Suggerimento}\normalsize\color{black}
\begin{lstlisting}
(defun fibonacci-fast (n)
  ;; a = F(i), b = F(i+1); a ogni passo a <- b, b <- a+b
  (labels ((iter (i a b)
             (if (= i 0)
                 a
                 (iter (1- i) b (+ a b)))))
    (iter n 0 1)))
\end{lstlisting}
\end{solbox}

\label{sol:55}
\begin{solbox}{55}
\small\color{black!55}\hyperref[ex:55]{$\leftarrow$ Esercizio}\enspace\textcolor{black!30}{|}\enspace\hyperref[hint:55]{$\leftarrow$ Suggerimento}\normalsize\color{black}
\begin{lstlisting}
(defun flatten* (tree)
  (cond
    ((null tree) nil)
    ;; atomo (foglia): avvolgilo in lista
    ((atom tree) (list tree))
    ;; nodo interno: appiattisci car e cdr separatamente
    (t (append (flatten* (car tree))
               (flatten* (cdr tree))))))
\end{lstlisting}
\end{solbox}

\label{sol:56}
\begin{solbox}{56}
\small\color{black!55}\hyperref[ex:56]{$\leftarrow$ Esercizio}\enspace\textcolor{black!30}{|}\enspace\hyperref[hint:56]{$\leftarrow$ Suggerimento}\normalsize\color{black}
\begin{lstlisting}
(defun tree-depth (tree)
  (cond
    ((null tree) 0)
    ((atom tree) 0)
    ;; profondita' = max(1 + profondita' car, profondita' cdr)
    (t (max (1+ (tree-depth (car tree)))
            (tree-depth (cdr tree))))))
\end{lstlisting}
\end{solbox}

\label{sol:57}
\begin{solbox}{57}
\small\color{black!55}\hyperref[ex:57]{$\leftarrow$ Esercizio}\normalsize\color{black}
\begin{lstlisting}
(defun count-atoms (tree)
  (cond ((null tree) 0)
        ((atom tree) 1)
        (t (+ (count-atoms (car tree))
              (count-atoms (cdr tree))))))
\end{lstlisting}
\end{solbox}

\label{sol:58}
\begin{solbox}{58}
\small\color{black!55}\hyperref[ex:58]{$\leftarrow$ Esercizio}\normalsize\color{black}
\begin{lstlisting}
(defun sum-tree (tree)
  (cond ((null tree)   0)
        ((numberp tree) tree)
        ((atom tree)    0)
        (t (+ (sum-tree (car tree))
              (sum-tree (cdr tree))))))
\end{lstlisting}
\end{solbox}

\label{sol:59}
\begin{solbox}{59}
\small\color{black!55}\hyperref[ex:59]{$\leftarrow$ Esercizio}\normalsize\color{black}
\begin{lstlisting}
(defun tree-member-p (item tree)
  (cond ((null tree) nil)
        ((equal tree item) t)
        ((atom tree) nil)
        (t (or (tree-member-p item (car tree))
               (tree-member-p item (cdr tree))))))
\end{lstlisting}
\end{solbox}

\label{sol:60}
\begin{solbox}{60}
\small\color{black!55}\hyperref[ex:60]{$\leftarrow$ Esercizio}\enspace\textcolor{black!30}{|}\enspace\hyperref[hint:60]{$\leftarrow$ Suggerimento}\normalsize\color{black}
\begin{lstlisting}
(defun sostituisci-profondo (old new tree)
  (cond ((null tree) nil)
        ((equal tree old) new)   ; sostituzione
        ((atom tree) tree)       ; atomo diverso: lascia invariato
        (t (cons (sostituisci-profondo old new (car tree))
                 (sostituisci-profondo old new (cdr tree))))))
\end{lstlisting}
\end{solbox}

\label{sol:61}
\begin{solbox}{61}
\small\color{black!55}\hyperref[ex:61]{$\leftarrow$ Esercizio}\enspace\textcolor{black!30}{|}\enspace\hyperref[hint:61]{$\leftarrow$ Suggerimento}\normalsize\color{black}
\begin{lstlisting}
(defun power-fast (base exp)
  (cond ((= exp 0) 1)
        ((evenp exp)
         ;; evita doppia ricorsione con let
         (let ((half (power-fast base (/ exp 2))))
           (* half half)))
        (t (* base (power-fast base (1- exp))))))
\end{lstlisting}
\end{solbox}

\label{sol:62}
\begin{solbox}{62}
\small\color{black!55}\hyperref[ex:62]{$\leftarrow$ Esercizio}\normalsize\color{black}
\begin{lstlisting}
(defun mcd* (a b)
  (if (= b 0) a (mcd* b (mod a b))))

(defun mcm* (a b)
  (/ (* a b) (mcd* a b)))
\end{lstlisting}
\end{solbox}

\label{sol:63}
\begin{solbox}{63}
\small\color{black!55}\hyperref[ex:63]{$\leftarrow$ Esercizio}\normalsize\color{black}
\begin{lstlisting}
(defun primo-p (n)
  (if (< n 2)
      nil
      (loop for d from 2 to (isqrt n)
            never (= (mod n d) 0))))
\end{lstlisting}
\end{solbox}

\label{sol:64}
\begin{solbox}{64}
\small\color{black!55}\hyperref[ex:64]{$\leftarrow$ Esercizio}\enspace\textcolor{black!30}{|}\enspace\hyperref[hint:64]{$\leftarrow$ Suggerimento}\normalsize\color{black}
\begin{lstlisting}
(defun crivello (n)
  ;; array di candidati: T = possibile primo
  (let ((sieve (make-array (1+ n) :initial-element t)))
    (setf (aref sieve 0) nil
          (aref sieve 1) nil)
    ;; elimina multipli di ogni primo trovato
    (loop for i from 2 to (isqrt n)
          when (aref sieve i)
          do (loop for j from (* i i) to n by i
                   do (setf (aref sieve j) nil)))
    ;; raccogli gli indici rimasti T
    (loop for i from 2 to n
          when (aref sieve i) collect i)))
\end{lstlisting}
\end{solbox}

\label{sol:65}
\begin{solbox}{65}
\small\color{black!55}\hyperref[ex:65]{$\leftarrow$ Esercizio}\enspace\textcolor{black!30}{|}\enspace\hyperref[hint:65]{$\leftarrow$ Suggerimento}\normalsize\color{black}
\begin{lstlisting}
(defun prime-factors (n)
  ;; prova divisori da 2 in su; non avanzare d se divide
  (labels ((iter (n d)
             (cond ((> (* d d) n)
                    (if (> n 1) (list n) nil))
                   ((= (mod n d) 0)
                    (cons d (iter (/ n d) d)))
                   (t (iter n (1+ d))))))
    (iter n 2)))
\end{lstlisting}
\end{solbox}

\label{sol:66}
\begin{solbox}{66}
\small\color{black!55}\hyperref[ex:66]{$\leftarrow$ Esercizio}\normalsize\color{black}
\begin{lstlisting}
(defun sublist (lst inizio fine)
  (let ((len (- fine inizio)))
    (if (<= len 0)
        nil
        (loop for i from inizio below fine
              collect (nth i lst)))))
\end{lstlisting}
\end{solbox}

\label{sol:67}
\begin{solbox}{67}
\small\color{black!55}\hyperref[ex:67]{$\leftarrow$ Esercizio}\enspace\textcolor{black!30}{|}\enspace\hyperref[hint:67]{$\leftarrow$ Suggerimento}\normalsize\color{black}
\begin{lstlisting}
(defun raggruppa (lst n)
  (if (null lst)
      nil
      (cons (loop for x in lst
                  for i below n collect x)
            (raggruppa (nthcdr n lst) n))))
\end{lstlisting}
\end{solbox}

\label{sol:68}
\begin{solbox}{68}
\small\color{black!55}\hyperref[ex:68]{$\leftarrow$ Esercizio}\normalsize\color{black}
\begin{lstlisting}
(defun interseca (lst1 lst2)
  (remove-if-not (lambda (x) (member x lst2)) lst1))
\end{lstlisting}
\end{solbox}

% =============================================================
\section*{Higher-order e closure}
% =============================================================

\label{sol:69}
\begin{solbox}{69}
\small\color{black!55}\hyperref[ex:69]{$\leftarrow$ Esercizio}\normalsize\color{black}
\begin{lstlisting}
(defun square-list (lst)
  (mapcar (lambda (x) (* x x)) lst))
\end{lstlisting}
\end{solbox}

\label{sol:70}
\begin{solbox}{70}
\small\color{black!55}\hyperref[ex:70]{$\leftarrow$ Esercizio}\normalsize\color{black}
\begin{lstlisting}
(defun even-list (lst)
  (remove-if-not #'evenp lst))
\end{lstlisting}
\end{solbox}

\label{sol:71}
\begin{solbox}{71}
\small\color{black!55}\hyperref[ex:71]{$\leftarrow$ Esercizio}\normalsize\color{black}
\begin{lstlisting}
(defun vocali-iniziali (lst)
  (remove-if-not
   (lambda (s)
     (member (char-downcase (char s 0))
             '(#\a #\e #\i #\o #\u)))
   lst))
\end{lstlisting}
\end{solbox}

\label{sol:72}
\begin{solbox}{72}
\small\color{black!55}\hyperref[ex:72]{$\leftarrow$ Esercizio}\normalsize\color{black}
\begin{lstlisting}
(defun sum-reduce (lst)
  (reduce #'+ lst :initial-value 0))
\end{lstlisting}
\end{solbox}

\label{sol:73}
\begin{solbox}{73}
\small\color{black!55}\hyperref[ex:73]{$\leftarrow$ Esercizio}\normalsize\color{black}
\begin{lstlisting}
(defun product-reduce (lst)
  (reduce #'* lst :initial-value 1))
\end{lstlisting}
\end{solbox}

\label{sol:74}
\begin{solbox}{74}
\small\color{black!55}\hyperref[ex:74]{$\leftarrow$ Esercizio}\normalsize\color{black}
\begin{lstlisting}
(defun reduce-min (lst)
  (if (null lst)
      (error "lista vuota")
      (reduce #'min lst)))

(defun reduce-max (lst)
  (if (null lst)
      (error "lista vuota")
      (reduce #'max lst)))
\end{lstlisting}
\end{solbox}

\label{sol:75}
\begin{solbox}{75}
\small\color{black!55}\hyperref[ex:75]{$\leftarrow$ Esercizio}\enspace\textcolor{black!30}{|}\enspace\hyperref[hint:75]{$\leftarrow$ Suggerimento}\normalsize\color{black}
\begin{lstlisting}
(defun cifre->intero (lst)
  ;; a ogni passo: acc = acc*10 + cifra
  (reduce (lambda (acc d) (+ (* acc 10) d))
          lst :initial-value 0))
\end{lstlisting}
\end{solbox}

\label{sol:76}
\begin{solbox}{76}
\small\color{black!55}\hyperref[ex:76]{$\leftarrow$ Esercizio}\normalsize\color{black}
\begin{lstlisting}
(defun appiattisci-un-livello (lst)
  (reduce #'append lst :initial-value nil))
\end{lstlisting}
\end{solbox}

\label{sol:77}
\begin{solbox}{77}
\small\color{black!55}\hyperref[ex:77]{$\leftarrow$ Esercizio}\enspace\textcolor{black!30}{|}\enspace\hyperref[hint:77]{$\leftarrow$ Suggerimento}\normalsize\color{black}
\begin{lstlisting}
(defun horner (coefficienti x)
  ;; fold: acc = acc*x + c
  (reduce (lambda (acc c) (+ (* acc x) c))
          coefficienti :initial-value 0))
\end{lstlisting}
\end{solbox}

\label{sol:78}
\begin{solbox}{78}
\small\color{black!55}\hyperref[ex:78]{$\leftarrow$ Esercizio}\normalsize\color{black}
\begin{lstlisting}
(defun prodotto-scalare (v1 v2)
  (apply #'+ (mapcar #'* v1 v2)))
\end{lstlisting}
\end{solbox}

\label{sol:79}
\begin{solbox}{79}
\small\color{black!55}\hyperref[ex:79]{$\leftarrow$ Esercizio}\normalsize\color{black}
\begin{lstlisting}
;; mapcan applica f a ogni elemento e concatena le liste risultanti;
;; mapcar le raccoglie senza appiattire.
(defun raddoppia (lst)
  (mapcan (lambda (x) (list x x)) lst))
\end{lstlisting}
\end{solbox}

\label{sol:80}
\begin{solbox}{80}
\small\color{black!55}\hyperref[ex:80]{$\leftarrow$ Esercizio}\normalsize\color{black}
\begin{lstlisting}
(defun ordina-per-punteggio (lst)
  (sort (copy-list lst)
        (lambda (a b) (> (cdr a) (cdr b)))))
\end{lstlisting}
\end{solbox}

\label{sol:81}
\begin{solbox}{81}
\small\color{black!55}\hyperref[ex:81]{$\leftarrow$ Esercizio}\normalsize\color{black}
\begin{lstlisting}
(defun tutti-p (pred lst)
  (cond ((null lst) t)
        ((funcall pred (car lst))
         (tutti-p pred (cdr lst)))
        (t nil)))

(defun qualcuno-p (pred lst)
  (cond ((null lst) nil)
        ((funcall pred (car lst)) t)
        (t (qualcuno-p pred (cdr lst)))))
\end{lstlisting}
\end{solbox}

\label{sol:82}
\begin{solbox}{82}
\small\color{black!55}\hyperref[ex:82]{$\leftarrow$ Esercizio}\normalsize\color{black}
\begin{lstlisting}
(defun count-matching (pred lst)
  (reduce (lambda (acc x)
            (if (funcall pred x) (1+ acc) acc))
          lst :initial-value 0))
\end{lstlisting}
\end{solbox}

\label{sol:83}
\begin{solbox}{83}
\small\color{black!55}\hyperref[ex:83]{$\leftarrow$ Esercizio}\enspace\textcolor{black!30}{|}\enspace\hyperref[hint:83]{$\leftarrow$ Suggerimento}\normalsize\color{black}
\begin{lstlisting}
(defun raggruppa-per (f lst)
  ;; per ogni elemento, trova o crea il gruppo nella alist
  (let ((result nil))
    (dolist (x lst result)
      (let* ((k (funcall f x))
             (bucket (assoc k result)))
        (if bucket
            (setf (cdr bucket) (append (cdr bucket) (list x)))
            (push (cons k (list x)) result))))))
\end{lstlisting}
\end{solbox}

\label{sol:84}
\begin{solbox}{84}
\small\color{black!55}\hyperref[ex:84]{$\leftarrow$ Esercizio}\normalsize\color{black}
\begin{lstlisting}
(defun applica-a-tutti (funzioni x)
  (mapcar (lambda (f) (funcall f x)) funzioni))
\end{lstlisting}
\end{solbox}

\label{sol:85}
\begin{solbox}{85}
\small\color{black!55}\hyperref[ex:85]{$\leftarrow$ Esercizio}\enspace\textcolor{black!30}{|}\enspace\hyperref[hint:85]{$\leftarrow$ Suggerimento}\normalsize\color{black}
\begin{lstlisting}
(defun compose (f g)
  ;; cattura f e g nella closure
  (lambda (x) (funcall f (funcall g x))))
\end{lstlisting}
\end{solbox}

\label{sol:86}
\begin{solbox}{86}
\small\color{black!55}\hyperref[ex:86]{$\leftarrow$ Esercizio}\enspace\textcolor{black!30}{|}\enspace\hyperref[hint:86]{$\leftarrow$ Suggerimento}\normalsize\color{black}
\begin{lstlisting}
(defun partial (funzione &rest argomenti)
  ;; cattura funzione e argomenti nella closure
  (lambda (&rest rest-args)
    (apply funzione (append argomenti rest-args))))
\end{lstlisting}
\end{solbox}

\label{sol:87}
\begin{solbox}{87}
\small\color{black!55}\hyperref[ex:87]{$\leftarrow$ Esercizio}\enspace\textcolor{black!30}{|}\enspace\hyperref[hint:87]{$\leftarrow$ Suggerimento}\normalsize\color{black}
\begin{lstlisting}
(defun make-counter (&optional (start 0) (step 1))
  ;; curr e' privato alla closure
  (let ((curr start))
    (lambda ()
      (prog1 curr       ; restituisce curr prima dell'incremento
        (incf curr step)))))
\end{lstlisting}
\end{solbox}

\label{sol:88}
\begin{solbox}{88}
\small\color{black!55}\hyperref[ex:88]{$\leftarrow$ Esercizio}\normalsize\color{black}
\begin{lstlisting}
(defun make-accumulator (iniziale)
  (let ((totale iniziale))
    (lambda (x)
      (incf totale x))))
\end{lstlisting}
\end{solbox}

\label{sol:89}
\begin{solbox}{89}
\small\color{black!55}\hyperref[ex:89]{$\leftarrow$ Esercizio}\enspace\textcolor{black!30}{|}\enspace\hyperref[hint:89]{$\leftarrow$ Suggerimento}\normalsize\color{black}
\begin{lstlisting}
(defun make-fib-generator ()
  ;; a e b condividono la stessa closure
  (let ((a 0) (b 1))
    (lambda ()
      (prog1 a
        (psetf a b           ; aggiorna simultaneamente
               b (+ a b))))))
\end{lstlisting}
\end{solbox}

\label{sol:90}
\begin{solbox}{90}
\small\color{black!55}\hyperref[ex:90]{$\leftarrow$ Esercizio}\normalsize\color{black}
\begin{lstlisting}
(defun pipeline (valore funzioni)
  (reduce (lambda (v f) (funcall f v))
          funzioni :initial-value valore))
\end{lstlisting}
\end{solbox}

\label{sol:91}
\begin{solbox}{91}
\small\color{black!55}\hyperref[ex:91]{$\leftarrow$ Esercizio}\enspace\textcolor{black!30}{|}\enspace\hyperref[hint:91]{$\leftarrow$ Suggerimento}\normalsize\color{black}
\begin{lstlisting}
(defun memoize (funzione)
  ;; la cache e' privata alla closure
  (let ((cache (make-hash-table :test #'equal)))
    (lambda (&rest args)
      (multiple-value-bind (val found)
          (gethash args cache)
        (if found
            val
            ;; calcola e memorizza
            (setf (gethash args cache)
                  (apply funzione args)))))))
\end{lstlisting}
\end{solbox}

\label{sol:92}
\begin{solbox}{92}
\small\color{black!55}\hyperref[ex:92]{$\leftarrow$ Esercizio}\normalsize\color{black}
\begin{lstlisting}
;; Mini libreria funzionale (estratto)
(defun map-list (f lst)
  "Applica F a ogni elemento di LST."
  (mapcar f lst))

(defun filter-list (pred lst)
  "Filtra LST mantenendo gli elementi per cui PRED e' vero."
  (remove-if-not pred lst))

(defun find-matching (pred lst)
  "Restituisce il primo elemento che soddisfa PRED."
  (find-if pred lst))

(defun take (n lst)
  "Restituisce i primi N elementi di LST."
  (if (or (= n 0) (null lst))
      nil
      (cons (car lst) (take (1- n) (cdr lst)))))

(defun drop (n lst)
  "Rimuove i primi N elementi di LST."
  (if (or (= n 0) (null lst))
      lst
      (drop (1- n) (cdr lst))))

(defun zip (lst1 lst2)
  "Combina due liste in lista di coppie."
  (if (or (null lst1) (null lst2))
      nil
      (cons (list (car lst1) (car lst2))
            (zip (cdr lst1) (cdr lst2)))))
\end{lstlisting}
\end{solbox}

% =============================================================
\section*{Strutture dati}
% =============================================================

\label{sol:93}
\begin{solbox}{93}
\small\color{black!55}\hyperref[ex:93]{$\leftarrow$ Esercizio}\normalsize\color{black}
\begin{lstlisting}
(defun lookup (key alist)
  (let ((pair (assoc key alist)))
    (if pair (cdr pair) nil)))

(defun set-key (key value alist)
  ;; costruisce nuova alist senza modificare l'originale
  (cons (cons key value)
        (remove-if (lambda (p) (equal (car p) key))
                   alist)))
\end{lstlisting}
\end{solbox}

\label{sol:94}
\begin{solbox}{94}
\small\color{black!55}\hyperref[ex:94]{$\leftarrow$ Esercizio}\enspace\textcolor{black!30}{|}\enspace\hyperref[hint:94]{$\leftarrow$ Suggerimento}\normalsize\color{black}
\begin{lstlisting}
(defun frequencies (lst)
  ;; per ogni elemento, aggiorna o aggiunge la coppia
  (let ((result nil))
    (dolist (x lst result)
      (let ((pair (assoc x result)))
        (if pair
            (incf (cdr pair))
            (push (cons x 1) result))))))
\end{lstlisting}
\end{solbox}

\label{sol:95}
\begin{solbox}{95}
\small\color{black!55}\hyperref[ex:95]{$\leftarrow$ Esercizio}\normalsize\color{black}
\begin{lstlisting}
(defun ht-get (key ht)
  (gethash key ht))

(defun ht-set (key value ht)
  (setf (gethash key ht) value))

(defun ht-delete (key ht)
  (remhash key ht))

(defun ht-keys (ht)
  (loop for k being the hash-keys of ht collect k))
\end{lstlisting}
\end{solbox}

\label{sol:96}
\begin{solbox}{96}
\small\color{black!55}\hyperref[ex:96]{$\leftarrow$ Esercizio}\normalsize\color{black}
\begin{lstlisting}
(defun frequencies-hash (lst)
  (let ((ht (make-hash-table)))
    (dolist (x lst ht)
      (incf (gethash x ht 0)))))
\end{lstlisting}
\end{solbox}

\label{sol:97}
\begin{solbox}{97}
\small\color{black!55}\hyperref[ex:97]{$\leftarrow$ Esercizio}\enspace\textcolor{black!30}{|}\enspace\hyperref[hint:97]{$\leftarrow$ Suggerimento}\normalsize\color{black}
\begin{lstlisting}
(defun frequenza-caratteri (stringa)
  (let ((ht (make-hash-table :test #'eql)))
    (loop for c across stringa
          do (incf (gethash c ht 0)))
    ht))
\end{lstlisting}
\end{solbox}

\label{sol:98}
\begin{solbox}{98}
\small\color{black!55}\hyperref[ex:98]{$\leftarrow$ Esercizio}\enspace\textcolor{black!30}{|}\enspace\hyperref[hint:98]{$\leftarrow$ Suggerimento}\normalsize\color{black}
\begin{lstlisting}
(defun unisci-hash (ht1 ht2)
  (let ((result (make-hash-table :test (hash-table-test ht1))))
    (maphash (lambda (k v) (setf (gethash k result) v)) ht1)
    (maphash (lambda (k v)
               (incf (gethash k result 0) v))
             ht2)
    result))
\end{lstlisting}
\end{solbox}

\label{sol:99}
\begin{solbox}{99}
\small\color{black!55}\hyperref[ex:99]{$\leftarrow$ Esercizio}\enspace\textcolor{black!30}{|}\enspace\hyperref[hint:99]{$\leftarrow$ Suggerimento}\normalsize\color{black}
\begin{lstlisting}
(defun inverti-hash (ht)
  (let ((result (make-hash-table :test #'equal)))
    (maphash (lambda (k v) (setf (gethash v result) k)) ht)
    result))
\end{lstlisting}
\end{solbox}

\label{sol:100}
\begin{solbox}{100}
\small\color{black!55}\hyperref[ex:100]{$\leftarrow$ Esercizio}\normalsize\color{black}
\begin{lstlisting}
(defun alist->hash (alist)
  (let ((ht (make-hash-table)))
    (dolist (pair alist ht)
      (setf (gethash (car pair) ht) (cdr pair)))))

(defun hash->alist (ht)
  (loop for k being the hash-keys of ht
        collect (cons k (gethash k ht))))
\end{lstlisting}
\end{solbox}

\label{sol:101}
\begin{solbox}{101}
\small\color{black!55}\hyperref[ex:101]{$\leftarrow$ Esercizio}\normalsize\color{black}
\begin{lstlisting}
(defun mat-get (m r c) (aref m r c))

(defun mat-set (m r c val) (setf (aref m r c) val))

(defun stampa-matrice (m)
  (let ((n (car (array-dimensions m))))
    (dotimes (r n)
      (dotimes (c n)
        (format t "~4D" (aref m r c)))
      (terpri))))
\end{lstlisting}
\end{solbox}

\label{sol:102}
\begin{solbox}{102}
\small\color{black!55}\hyperref[ex:102]{$\leftarrow$ Esercizio}\enspace\textcolor{black!30}{|}\enspace\hyperref[hint:102]{$\leftarrow$ Suggerimento}\normalsize\color{black}
\begin{lstlisting}
(defun trasposta (m n)
  ;; t[i][j] = m[j][i]
  (let ((result (make-array (list n n))))
    (dotimes (i n result)
      (dotimes (j n)
        (setf (aref result i j)
              (aref m j i))))))
\end{lstlisting}
\end{solbox}

\label{sol:103}
\begin{solbox}{103}
\small\color{black!55}\hyperref[ex:103]{$\leftarrow$ Esercizio}\normalsize\color{black}
\begin{lstlisting}
(defun inverti-array! (v)
  ;; scambia elementi simmetrici senza lista ausiliaria
  (let ((len (length v)))
    (dotimes (i (floor len 2))
      (let ((tmp (aref v i)))
        (setf (aref v i) (aref v (- len 1 i))
              (aref v (- len 1 i)) tmp)))))
\end{lstlisting}
\end{solbox}

\label{sol:104}
\begin{solbox}{104}
\small\color{black!55}\hyperref[ex:104]{$\leftarrow$ Esercizio}\enspace\textcolor{black!30}{|}\enspace\hyperref[hint:104]{$\leftarrow$ Suggerimento}\normalsize\color{black}
\begin{lstlisting}
(defun moltiplica-matrici (a b n)
  (let ((c (make-array (list n n) :initial-element 0)))
    (dotimes (i n c)
      (dotimes (j n)
        (dotimes (k n)
          (incf (aref c i j)
                (* (aref a i k) (aref b k j))))))))
\end{lstlisting}
\end{solbox}

\label{sol:105}
\begin{solbox}{105}
\small\color{black!55}\hyperref[ex:105]{$\leftarrow$ Esercizio}\normalsize\color{black}
\begin{lstlisting}
(defstruct persona
  nome
  (eta 18)
  citta)

(defun crea-persona (nome eta citta)
  (when (not (and (integerp eta) (> eta 0)))
    (error "Eta' non valida: ~A" eta))
  (make-persona :nome nome :eta eta :citta citta))
\end{lstlisting}
\end{solbox}

\label{sol:106}
\begin{solbox}{106}
\small\color{black!55}\hyperref[ex:106]{$\leftarrow$ Esercizio}\normalsize\color{black}
\begin{lstlisting}
(defstruct nodo
  valore
  (sinistra nil)
  (destra   nil))

;; albero:     5
;;            / \
;;           3   8
;;          / \
;;         1   4
(defparameter *albero*
  (make-nodo :valore 5
    :sinistra (make-nodo :valore 3
                :sinistra (make-nodo :valore 1)
                :destra   (make-nodo :valore 4))
    :destra   (make-nodo :valore 8)))
\end{lstlisting}
\end{solbox}

\label{sol:107}
\begin{solbox}{107}
\small\color{black!55}\hyperref[ex:107]{$\leftarrow$ Esercizio}\normalsize\color{black}
\begin{lstlisting}
;; Rappresentazione: (valore sinistra destra)
;; NIL = albero vuoto
(defun make-leaf (v) (list v nil nil))
(defun make-node (v l r) (list v l r))
(defun node-val  (n) (car n))
(defun node-left (n) (cadr n))
(defun node-right (n) (caddr n))
\end{lstlisting}
\end{solbox}

\label{sol:108}
\begin{solbox}{108}
\small\color{black!55}\hyperref[ex:108]{$\leftarrow$ Esercizio}\enspace\textcolor{black!30}{|}\enspace\hyperref[hint:108]{$\leftarrow$ Suggerimento}\normalsize\color{black}
\begin{lstlisting}
(defun bst-insert (value tree)
  ;; restituisce il nuovo albero (non modifica in loco)
  (if (null tree)
      (make-nodo :valore value)
      (cond ((< value (nodo-valore tree))
             (make-nodo
               :valore (nodo-valore tree)
               :sinistra (bst-insert value (nodo-sinistra tree))
               :destra   (nodo-destra tree)))
            ((> value (nodo-valore tree))
             (make-nodo
               :valore (nodo-valore tree)
               :sinistra (nodo-sinistra tree)
               :destra   (bst-insert value (nodo-destra tree))))
            ;; duplicato: non inserire
            (t tree))))

(defun bst-find (value tree)
  (cond ((null tree) nil)
        ((= value (nodo-valore tree)) t)
        ((< value (nodo-valore tree))
         (bst-find value (nodo-sinistra tree)))
        (t (bst-find value (nodo-destra tree)))))
\end{lstlisting}
\end{solbox}

\label{sol:109}
\begin{solbox}{109}
\small\color{black!55}\hyperref[ex:109]{$\leftarrow$ Esercizio}\enspace\textcolor{black!30}{|}\enspace\hyperref[hint:109]{$\leftarrow$ Suggerimento}\normalsize\color{black}
\begin{lstlisting}
(defun preorder (tree)
  (if (null tree) nil
      (append (list (nodo-valore tree))
              (preorder (nodo-sinistra tree))
              (preorder (nodo-destra tree)))))

(defun inorder (tree)
  (if (null tree) nil
      (append (inorder (nodo-sinistra tree))
              (list (nodo-valore tree))
              (inorder (nodo-destra tree)))))

(defun postorder (tree)
  (if (null tree) nil
      (append (postorder (nodo-sinistra tree))
              (postorder (nodo-destra tree))
              (list (nodo-valore tree)))))
\end{lstlisting}
\end{solbox}

\label{sol:110}
\begin{solbox}{110}
\small\color{black!55}\hyperref[ex:110]{$\leftarrow$ Esercizio}\normalsize\color{black}
\begin{lstlisting}
(defun tree-height (tree)
  (if (null tree) 0
      (1+ (max (tree-height (nodo-sinistra tree))
               (tree-height (nodo-destra tree))))))

(defun count-leaves (tree)
  (cond ((null tree) 0)
        ((and (null (nodo-sinistra tree))
              (null (nodo-destra tree))) 1)
        (t (+ (count-leaves (nodo-sinistra tree))
              (count-leaves (nodo-destra tree))))))
\end{lstlisting}
\end{solbox}

\label{sol:111}
\begin{solbox}{111}
\small\color{black!55}\hyperref[ex:111]{$\leftarrow$ Esercizio}\normalsize\color{black}
\begin{lstlisting}
(defun specchia (tree)
  (if (null tree) nil
      (make-nodo :valore   (nodo-valore tree)
                 :sinistra (specchia (nodo-destra tree))
                 :destra   (specchia (nodo-sinistra tree)))))
\end{lstlisting}
\end{solbox}

\label{sol:112}
\begin{solbox}{112}
\small\color{black!55}\hyperref[ex:112]{$\leftarrow$ Esercizio}\enspace\textcolor{black!30}{|}\enspace\hyperref[hint:112]{$\leftarrow$ Suggerimento}\normalsize\color{black}
\begin{lstlisting}
(defun bilanciato-p (tree)
  ;; ritorna altezza se bilanciato, -1 come segnale di errore
  (labels ((check (t)
             (if (null t)
                 0
                 (let ((hl (check (nodo-sinistra t)))
                       (hr (check (nodo-destra t))))
                   (if (or (= hl -1) (= hr -1)
                           (> (abs (- hl hr)) 1))
                       -1
                       (1+ (max hl hr)))))))
    (/= (check tree) -1)))
\end{lstlisting}
\end{solbox}

\label{sol:113}
\begin{solbox}{113}
\small\color{black!55}\hyperref[ex:113]{$\leftarrow$ Esercizio}\normalsize\color{black}
\begin{lstlisting}
;; Pila come lista avvolta in cons-cell per mutabilita'
(defun make-stack () (list nil))

(defun stack-push (item stack)
  (setf (car stack) (cons item (car stack))))

(defun stack-pop (stack)
  (prog1 (caar stack)
    (setf (car stack) (cdar stack))))

(defun stack-peek (stack) (caar stack))

(defun stack-empty-p (stack) (null (car stack)))
\end{lstlisting}
\end{solbox}

\label{sol:114}
\begin{solbox}{114}
\small\color{black!55}\hyperref[ex:114]{$\leftarrow$ Esercizio}\enspace\textcolor{black!30}{|}\enspace\hyperref[hint:114]{$\leftarrow$ Suggerimento}\normalsize\color{black}
\begin{lstlisting}
;; Coda con due liste: testa (estrai) e fondo (aggiungi)
(defstruct queue (head nil) (tail nil))

(defun enqueue (item q)
  (setf (queue-tail q) (cons item (queue-tail q))))

(defun dequeue (q)
  ;; se testa vuota, ribalta il fondo
  (when (null (queue-head q))
    (setf (queue-head q) (reverse (queue-tail q))
          (queue-tail q) nil))
  (prog1 (car (queue-head q))
    (setf (queue-head q) (cdr (queue-head q)))))

(defun queue-empty-p (q)
  (and (null (queue-head q))
       (null (queue-tail q))))
\end{lstlisting}
\end{solbox}

\label{sol:115}
\begin{solbox}{115}
\small\color{black!55}\hyperref[ex:115]{$\leftarrow$ Esercizio}\enspace\textcolor{black!30}{|}\enspace\hyperref[hint:115]{$\leftarrow$ Suggerimento}\normalsize\color{black}
\begin{lstlisting}
(defun valida-parentesi (stringa)
  ;; mappa ogni chiusa alla corrispondente aperta
  (let ((match '((#\) . #\()
                 (#\] . #\[)
                 (#\} . #\{)))
        (stack (make-stack)))
    (loop for c across stringa
          do (cond
               ;; aperta: push
               ((member c '(#\( #\[ #\{))
                (stack-push c stack))
               ;; chiusa: verifica top
               ((assoc c match)
                (if (or (stack-empty-p stack)
                        (not (char= (stack-pop stack)
                                    (cdr (assoc c match)))))
                    (return-from valida-parentesi nil)))))
    (stack-empty-p stack)))
\end{lstlisting}
\end{solbox}

\label{sol:116}
\begin{solbox}{116}
\small\color{black!55}\hyperref[ex:116]{$\leftarrow$ Esercizio}\normalsize\color{black}
\begin{lstlisting}
(defparameter *rubrica* (make-hash-table :test #'equal))

(defun add-contact (name phone)
  (setf (gethash name *rubrica*) phone))

(defun find-contact (name)
  (or (gethash name *rubrica*)
      (error "Contatto non trovato: ~A" name)))

(defun remove-contact (name)
  (remhash name *rubrica*))

(defun list-contacts ()
  (maphash (lambda (n p) (format t "~A: ~A~%" n p))
           *rubrica*))
\end{lstlisting}
\end{solbox}

\label{sol:117}
\begin{solbox}{117}
\small\color{black!55}\hyperref[ex:117]{$\leftarrow$ Esercizio}\normalsize\color{black}
\begin{lstlisting}
(defun neighbors (node graph)
  (cdr (assoc node graph)))

(defun add-edge (n1 n2 graph)
  ;; aggiunge arco bidirezionale
  (let ((g (copy-tree graph)))
    (let ((entry1 (assoc n1 g))
          (entry2 (assoc n2 g)))
      (when entry1 (setf (cdr entry1) (cons n2 (cdr entry1))))
      (when entry2 (setf (cdr entry2) (cons n1 (cdr entry2))))
      g)))
\end{lstlisting}
\end{solbox}

\label{sol:118}
\begin{solbox}{118}
\small\color{black!55}\hyperref[ex:118]{$\leftarrow$ Esercizio}\enspace\textcolor{black!30}{|}\enspace\hyperref[hint:118]{$\leftarrow$ Suggerimento}\normalsize\color{black}
\begin{lstlisting}
(defun dfs (graph start)
  ;; visited e' condivisa dalla closure tramite let
  (let ((visited nil))
    (labels ((visit (node)
               (unless (member node visited)
                 (push node visited)
                 (dolist (n (neighbors node graph))
                   (visit n)))))
      (visit start)
      (reverse visited))))
\end{lstlisting}
\end{solbox}

\label{sol:119}
\begin{solbox}{119}
\small\color{black!55}\hyperref[ex:119]{$\leftarrow$ Esercizio}\enspace\textcolor{black!30}{|}\enspace\hyperref[hint:119]{$\leftarrow$ Suggerimento}\normalsize\color{black}
\begin{lstlisting}
(defun bfs (graph start)
  (let ((visited nil)
        (queue   (list start)))
    (loop while queue
          for node = (pop queue)
          unless (member node visited)
          do (push node visited)
             (setf queue
                   (append queue
                           (remove-if (lambda (n) (member n visited))
                                      (neighbors node graph)))))
    (reverse visited)))
\end{lstlisting}
\end{solbox}

\label{sol:120}
\begin{solbox}{120}
\small\color{black!55}\hyperref[ex:120]{$\leftarrow$ Esercizio}\enspace\textcolor{black!30}{|}\enspace\hyperref[hint:120]{$\leftarrow$ Suggerimento}\normalsize\color{black}
\begin{lstlisting}
(defun tokenize (testo)
  ;; divide per spazi
  (let ((parole nil) (corrente ""))
    (loop for c across testo
          do (if (char= c #\Space)
                 (unless (string= corrente "")
                   (push corrente parole)
                   (setf corrente ""))
                 (setf corrente (concatenate 'string corrente (string c)))))
    (unless (string= corrente "")
      (push corrente parole))
    (reverse parole)))

(defun conta-parole (testo)
  (let ((ht (make-hash-table :test #'equal)))
    (dolist (w (tokenize testo) ht)
      (incf (gethash w ht 0)))))

(defun parola-piu-frequente (testo)
  (let ((ht (conta-parole testo))
        (best-w nil) (best-n 0))
    (maphash (lambda (w n)
               (when (> n best-n)
                 (setf best-w w best-n n)))
             ht)
    best-w))
\end{lstlisting}
\end{solbox}

% =============================================================
\section*{Macro}
% =============================================================

\label{sol:121}
\begin{solbox}{121}
\small\color{black!55}\hyperref[ex:121]{$\leftarrow$ Esercizio}\enspace\textcolor{black!30}{|}\enspace\hyperref[hint:121]{$\leftarrow$ Suggerimento}\normalsize\color{black}
\begin{lstlisting}
(defmacro quando (test &body corpo)
  ;; backquote: ,test inserisce il valore di test
  ;; ,@corpo: splice della lista corpo nel progn
  `(if ,test
       (progn ,@corpo)
       nil))

;; Verifica espansione:
;; (macroexpand-1 '(quando (> x 0) (print "ok") x))
;; -> (IF (> X 0) (PROGN (PRINT "ok") X) NIL)
\end{lstlisting}
\end{solbox}

\label{sol:122}
\begin{solbox}{122}
\small\color{black!55}\hyperref[ex:122]{$\leftarrow$ Esercizio}\enspace\textcolor{black!30}{|}\enspace\hyperref[hint:122]{$\leftarrow$ Suggerimento}\normalsize\color{black}
\begin{lstlisting}
(defmacro mio-while (test &body corpo)
  ;; usa loop nell'espansione: semplice ed idiomatico
  `(loop while ,test do ,@corpo))

;; Alternativa con tagbody/go (piu' primitiva):
;; (defmacro mio-while (test &body corpo)
;;   (let ((inizio (gensym)) (fine (gensym)))
;;     `(tagbody
;;        ,inizio
;;        (unless ,test (go ,fine))
;;        ,@corpo
;;        (go ,inizio)
;;        ,fine)))
\end{lstlisting}
\end{solbox}

\label{sol:123}
\begin{solbox}{123}
\small\color{black!55}\hyperref[ex:123]{$\leftarrow$ Esercizio}\enspace\textcolor{black!30}{|}\enspace\hyperref[hint:123]{$\leftarrow$ Suggerimento}\normalsize\color{black}
\begin{lstlisting}
(defmacro mio-for ((var da a) &body corpo)
  ;; gensym per start/end: evita valutazioni multiple
  ;; di da e a (effetti collaterali)
  (let ((start (gensym "START"))
        (end   (gensym "END")))
    `(let ((,start ,da)
           (,end   ,a))
       (loop for ,var from ,start to ,end
             do ,@corpo))))
\end{lstlisting}
\end{solbox}

% =============================================================
\section*{Backtracking e algoritmi}
% =============================================================

\label{sol:124}
\begin{solbox}{124}
\small\color{black!55}\hyperref[ex:124]{$\leftarrow$ Esercizio}\enspace\textcolor{black!30}{|}\enspace\hyperref[hint:124]{$\leftarrow$ Suggerimento}\normalsize\color{black}
\begin{lstlisting}
(defun subsets (lst)
  ;; ogni elemento: incluso o non incluso
  (if (null lst)
      '(())
      (let ((rest (subsets (cdr lst))))
        (append rest
                (mapcar (lambda (s)
                          (cons (car lst) s))
                        rest)))))
\end{lstlisting}
\end{solbox}

\label{sol:125}
\begin{solbox}{125}
\small\color{black!55}\hyperref[ex:125]{$\leftarrow$ Esercizio}\enspace\textcolor{black!30}{|}\enspace\hyperref[hint:125]{$\leftarrow$ Suggerimento}\normalsize\color{black}
\begin{lstlisting}
(defun permutations (lst)
  (if (null lst)
      '(())
      ;; per ogni elemento x: costruisci le permutazioni del resto
      (mapcan (lambda (x)
                (let ((rest (remove x lst :count 1)))
                  (mapcar (lambda (p) (cons x p))
                          (permutations rest))))
              lst)))
\end{lstlisting}
\end{solbox}

\label{sol:126}
\begin{solbox}{126}
\small\color{black!55}\hyperref[ex:126]{$\leftarrow$ Esercizio}\enspace\textcolor{black!30}{|}\enspace\hyperref[hint:126]{$\leftarrow$ Suggerimento}\normalsize\color{black}
\begin{lstlisting}
(defun combinations (lst k)
  (cond ((= k 0)    '(()))
        ((null lst)  nil)
        ;; con (car lst): combina k-1 dal resto; senza: combina k dal resto
        (t (append
            (mapcar (lambda (c) (cons (car lst) c))
                    (combinations (cdr lst) (1- k)))
            (combinations (cdr lst) k)))))
\end{lstlisting}
\end{solbox}

\label{sol:127}
\begin{solbox}{127}
\small\color{black!55}\hyperref[ex:127]{$\leftarrow$ Esercizio}\enspace\textcolor{black!30}{|}\enspace\hyperref[hint:127]{$\leftarrow$ Suggerimento}\normalsize\color{black}
\begin{lstlisting}
(defun somma-a-target (numeri target)
  (cond ((= target 0) '(()))
        ((or (null numeri) (< target 0)) nil)
        (t (let ((x (car numeri))
                 (rest (cdr numeri)))
             ;; includi x
             (append
              (mapcar (lambda (s) (cons x s))
                      (somma-a-target rest (- target x)))
              ;; non includere x
              (somma-a-target rest target))))))
\end{lstlisting}
\end{solbox}

\label{sol:128}
\begin{solbox}{128}
\small\color{black!55}\hyperref[ex:128]{$\leftarrow$ Esercizio}\enspace\textcolor{black!30}{|}\enspace\hyperref[hint:128]{$\leftarrow$ Suggerimento}\normalsize\color{black}
\begin{lstlisting}
(defun genera-parentesi (n)
  (let ((risultati nil))
    (labels ((gen (s aperte chiuse)
               (cond
                 ((= (length s) (* 2 n))
                  (push s risultati))
                 (t
                  (when (< aperte n)
                    (gen (concatenate 'string s "(")
                         (1+ aperte) chiuse))
                  (when (< chiuse aperte)
                    (gen (concatenate 'string s ")")
                         aperte (1+ chiuse)))))))
      (gen "" 0 0)
      (reverse risultati))))
\end{lstlisting}
\end{solbox}

\label{sol:129}
\begin{solbox}{129}
\small\color{black!55}\hyperref[ex:129]{$\leftarrow$ Esercizio}\enspace\textcolor{black!30}{|}\enspace\hyperref[hint:129]{$\leftarrow$ Suggerimento}\normalsize\color{black}
\begin{lstlisting}
(defun merge-sorted (lst1 lst2)
  (cond ((null lst1) lst2)
        ((null lst2) lst1)
        ((< (car lst1) (car lst2))
         (cons (car lst1) (merge-sorted (cdr lst1) lst2)))
        (t
         (cons (car lst2) (merge-sorted lst1 (cdr lst2))))))

(defun merge-sort* (lst)
  (let ((len (length lst)))
    (if (< len 2)
        lst
        (let* ((mid  (floor len 2))
               (left  (subseq lst 0 mid))
               (right (subseq lst mid)))
          (merge-sorted (merge-sort* left)
                        (merge-sort* right))))))
\end{lstlisting}
\end{solbox}

\label{sol:130}
\begin{solbox}{130}
\small\color{black!55}\hyperref[ex:130]{$\leftarrow$ Esercizio}\enspace\textcolor{black!30}{|}\enspace\hyperref[hint:130]{$\leftarrow$ Suggerimento}\normalsize\color{black}
\begin{lstlisting}
(defun quicksort* (lst)
  (if (or (null lst) (null (cdr lst)))
      lst
      (let* ((pivot   (car lst))
             (resto   (cdr lst))
             (minori  (remove-if-not (lambda (x) (< x pivot)) resto))
             (maggiori (remove-if (lambda (x) (< x pivot)) resto)))
        (append (quicksort* minori)
                (list pivot)
                (quicksort* maggiori)))))
\end{lstlisting}
\end{solbox}

\label{sol:131}
\begin{solbox}{131}
\small\color{black!55}\hyperref[ex:131]{$\leftarrow$ Esercizio}\enspace\textcolor{black!30}{|}\enspace\hyperref[hint:131]{$\leftarrow$ Suggerimento}\normalsize\color{black}
\begin{lstlisting}
(defun hanoi (n da a tramite)
  (when (> n 0)
    ;; passo 1: sposta n-1 da DA a TRAMITE (usando A)
    (hanoi (1- n) da tramite a)
    ;; passo 2: sposta il disco n da DA ad A
    (format t "Sposta disco ~D: ~A -> ~A~%" n da a)
    ;; passo 3: sposta n-1 da TRAMITE ad A (usando DA)
    (hanoi (1- n) tramite a da)))
\end{lstlisting}
\end{solbox}

\label{sol:132}
\begin{solbox}{132}
\small\color{black!55}\hyperref[ex:132]{$\leftarrow$ Esercizio}\enspace\textcolor{black!30}{|}\enspace\hyperref[hint:132]{$\leftarrow$ Suggerimento}\normalsize\color{black}
\begin{ragionamentoBox}{132}
La soluzione ricorsiva esplora tutte le $2^n$ combinazioni.
Per ogni oggetto ci sono due scelte: prenderlo o lasciarlo.
Restituiamo il valore massimo tra i due rami.
Complessita': $O(2^n)$ — accettabile per la versione didattica.
Per prestazioni migliori si userebbe la programmazione
dinamica con una tabella $n \times W$.
\end{ragionamentoBox}
\begin{lstlisting}
(defun knapsack (oggetti capacita)
  ;; oggetti: lista di (peso valore)
  (if (or (null oggetti) (= capacita 0))
      0
      (let* ((obj   (car oggetti))
             (resto (cdr oggetti))
             (peso  (car obj))
             (val   (cadr obj))
             ;; senza questo oggetto
             (senza (knapsack resto capacita)))
        (if (> peso capacita)
            senza
            ;; massimo tra prendere e non prendere
            (max senza
                 (+ val (knapsack resto (- capacita peso))))))))
\end{lstlisting}
\end{solbox}

\label{sol:133}
\begin{solbox}{133}
\small\color{black!55}\hyperref[ex:133]{$\leftarrow$ Esercizio}\enspace\textcolor{black!30}{|}\enspace\hyperref[hint:133]{$\leftarrow$ Suggerimento}\normalsize\color{black}
\begin{ragionamentoBox}{133}
DFS con backtracking su griglia. Usiamo un array separato
per i visitati per non modificare il labirinto originale.
La funzione restituisce il percorso se trovato, \fn{NIL} altrimenti.
\end{ragionamentoBox}
\begin{lstlisting}
(defun find-path (maze start goal)
  (let* ((rows (array-dimension maze 0))
         (cols (array-dimension maze 1))
         (visited (make-array (list rows cols) :initial-element nil)))
    (labels
      ((dfs (r c path)
         (cond
           ;; fuori dai bordi o muro o gia' visitata
           ((or (< r 0) (>= r rows) (< c 0) (>= c cols)
                (= (aref maze r c) 1)
                (aref visited r c))
            nil)
           ;; goal raggiunto
           ((equal (list r c) goal)
            (reverse (cons (list r c) path)))
           (t
            (setf (aref visited r c) t)
            ;; prova le 4 direzioni
            (let ((result (or (dfs (1+ r) c (cons (list r c) path))
                              (dfs (1- r) c (cons (list r c) path))
                              (dfs r (1+ c) (cons (list r c) path))
                              (dfs r (1- c) (cons (list r c) path)))))
              (setf (aref visited r c) nil) ; backtrack
              result)))))
      (dfs (car start) (cadr start) nil))))
\end{lstlisting}
\end{solbox}

\label{sol:134}
\begin{solbox}{134}
\small\color{black!55}\hyperref[ex:134]{$\leftarrow$ Esercizio}\enspace\textcolor{black!30}{|}\enspace\hyperref[hint:134]{$\leftarrow$ Suggerimento}\normalsize\color{black}
\begin{ragionamentoBox}{134}
\textbf{Rappresentazione:} vettore \fn{board} di \fn{n} interi, dove
\fn{(aref board col)} e' la riga della regina nella colonna \fn{col}.
Posizioniamo una regina per colonna, eliminando subito i conflitti.

\textbf{Pruning:} per la colonna \fn{col} e la riga \fn{row},
verifichiamo che nessuna colonna precedente \fn{pc} abbia:
\begin{itemize}[leftmargin=2em]
  \item stessa riga: \fn{(= (aref board pc) row)};
  \item stessa diagonale: \fn{(= (abs (- row (aref board pc))) (- col pc))}.
\end{itemize}
\textbf{Complessita':} con pruning, molto meglio di $O(n^n)$;
per $n=8$ si trovano 92 soluzioni in poche ms.
\end{ragionamentoBox}
\begin{lstlisting}
(defun n-queens (n)
  (let ((board (make-array n :initial-element 0))
        (solutions nil))
    (labels
      ((safe-p (col row)
         ;; nessun conflitto con le regine gia' posizionate
         (loop for pc below col
               never (let ((pr (aref board pc)))
                       (or (= pr row)
                           (= (abs (- row pr))
                              (- col pc))))))
       (place (col)
         (if (= col n)
             ;; soluzione trovata: registra una copia
             (push (coerce board 'list) solutions)
             (dotimes (row n)
               (when (safe-p col row)
                 (setf (aref board col) row)
                 (place (1+ col)))))))
      (place 0)
      solutions)))
\end{lstlisting}
\end{solbox}

\label{sol:135}
\begin{solbox}{135}
\small\color{black!55}\hyperref[ex:135]{$\leftarrow$ Esercizio}\enspace\textcolor{black!30}{|}\enspace\hyperref[hint:135]{$\leftarrow$ Suggerimento}\normalsize\color{black}
\solriferimento
\begin{ragionamentoBox}{135}
Backtracking con tre array di flag per righe, colonne e blocchi
$3\times3$ per un check $O(1)$ dei vincoli.
\end{ragionamentoBox}
\begin{lstlisting}
(defun solve-sudoku (board)
  ;; trova prima cella vuota
  (labels
    ((find-empty ()
       (loop for r below 9 do
         (loop for c below 9
               when (= (aref board r c) 0)
               do (return-from find-empty (list r c))))
       nil)
     (valid-p (r c val)
       ;; controlla riga, colonna, blocco 3x3
       (and (loop for i below 9 always (/= (aref board r i) val))
            (loop for i below 9 always (/= (aref board i c) val))
            (let ((br (* 3 (floor r 3)))
                  (bc (* 3 (floor c 3))))
              (loop for dr below 3 always
                (loop for dc below 3 always
                  (/= (aref board (+ br dr) (+ bc dc)) val)))))))
    (let ((empty (find-empty)))
      (if (null empty)
          board   ; risolto
          (let ((r (car empty)) (c (cadr empty)))
            (loop for val from 1 to 9
                  when (valid-p r c val)
                  do (setf (aref board r c) val)
                     (let ((result (solve-sudoku board)))
                       (when result (return-from solve-sudoku result)))
                     (setf (aref board r c) 0))
            nil))))) ; backtrack
\end{lstlisting}
\end{solbox}

\label{sol:136}
\begin{solbox}{136}
\small\color{black!55}\hyperref[ex:136]{$\leftarrow$ Esercizio}\enspace\textcolor{black!30}{|}\enspace\hyperref[hint:136]{$\leftarrow$ Suggerimento}\normalsize\color{black}
\solriferimento
\begin{ragionamentoBox}{136}
Backtracking con euristica di Warnsdorff: tra le mosse
disponibili, scegliamo sempre quella che porta alla cella
con meno uscite (mosse future disponibili).
Riduce drasticamente il backtracking rispetto al DFS puro.
\end{ragionamentoBox}
\begin{lstlisting}
(defparameter *moves*
  '((2 1)(2 -1)(-2 1)(-2 -1)(1 2)(1 -2)(-1 2)(-1 -2)))

(defun knights-tour (n)
  (let ((board (make-array (list n n) :initial-element 0)))
    (labels
      ((valid (r c)
         (and (>= r 0) (< r n) (>= c 0) (< c n)
              (= (aref board r c) 0)))
       (degree (r c)
         (count-if (lambda (m)
                     (valid (+ r (car m)) (+ c (cadr m))))
                   *moves*))
       (solve (r c step)
         (setf (aref board r c) step)
         (if (= step (* n n))
             t
             (let* ((next (remove-if-not
                           (lambda (m)
                             (valid (+ r (car m)) (+ c (cadr m))))
                           *moves*))
                    (sorted (sort next #'<
                                  :key (lambda (m)
                                         (degree (+ r (car m))
                                                 (+ c (cadr m)))))))
               (or (loop for m in sorted
                         thereis (solve (+ r (car m))
                                        (+ c (cadr m))
                                        (1+ step)))
                   (progn (setf (aref board r c) 0) nil))))))
      (when (solve 0 0 1) board))))
\end{lstlisting}
\end{solbox}

% =============================================================
\section*{Stringhe e I/O}
% =============================================================

\label{sol:137}
\begin{solbox}{137}
\small\color{black!55}\hyperref[ex:137]{$\leftarrow$ Esercizio}\normalsize\color{black}
\begin{lstlisting}
(defun reverse-string (s)
  (coerce (reverse (coerce s 'list)) 'string))
\end{lstlisting}
\end{solbox}

\label{sol:138}
\begin{solbox}{138}
\small\color{black!55}\hyperref[ex:138]{$\leftarrow$ Esercizio}\normalsize\color{black}
\begin{lstlisting}
(defun palindrome-p (s)
  (equal s (reverse-string s)))
\end{lstlisting}
\end{solbox}

\label{sol:139}
\begin{solbox}{139}
\small\color{black!55}\hyperref[ex:139]{$\leftarrow$ Esercizio}\enspace\textcolor{black!30}{|}\enspace\hyperref[hint:139]{$\leftarrow$ Suggerimento}\normalsize\color{black}
\begin{lstlisting}
(defun palindrome-robust-p (s)
  ;; normalizza: minuscolo, solo caratteri alfabetici
  (let ((clean (coerce
                (remove-if-not #'alpha-char-p
                               (coerce (string-downcase s) 'list))
                'string)))
    (equal clean (reverse-string clean))))
\end{lstlisting}
\end{solbox}

\label{sol:140}
\begin{solbox}{140}
\small\color{black!55}\hyperref[ex:140]{$\leftarrow$ Esercizio}\enspace\textcolor{black!30}{|}\enspace\hyperref[hint:140]{$\leftarrow$ Suggerimento}\normalsize\color{black}
\begin{lstlisting}
(defun anagram-p (s1 s2)
  ;; ordina le lettere di entrambe e confronta
  (equal (sort (coerce (string-downcase s1) 'list) #'char<)
         (sort (coerce (string-downcase s2) 'list) #'char<)))
\end{lstlisting}
\end{solbox}

\label{sol:141}
\begin{solbox}{141}
\small\color{black!55}\hyperref[ex:141]{$\leftarrow$ Esercizio}\normalsize\color{black}
\begin{lstlisting}
(defun count-char (c s)
  (count c s))
\end{lstlisting}
\end{solbox}

% =============================================================
\section*{Parsing e interprete}
% =============================================================

\label{sol:142}
\begin{solbox}{142}
\small\color{black!55}\hyperref[ex:142]{$\leftarrow$ Esercizio}\enspace\textcolor{black!30}{|}\enspace\hyperref[hint:142]{$\leftarrow$ Suggerimento}\normalsize\color{black}
\begin{lstlisting}
(defun rpn-eval (expr)
  (let ((stack (make-stack)))
    (dolist (token expr (stack-pop stack))
      (case token
        ((+ - * /)
         (let* ((b (stack-pop stack))
                (a (stack-pop stack)))
           (stack-push
            (case token
              (+ (+ a b)) (- (- a b))
              (* (* a b)) (/ (/ a b)))
            stack)))
        (otherwise
         (stack-push token stack))))))
\end{lstlisting}
\end{solbox}

\label{sol:143}
\begin{solbox}{143}
\small\color{black!55}\hyperref[ex:143]{$\leftarrow$ Esercizio}\enspace\textcolor{black!30}{|}\enspace\hyperref[hint:143]{$\leftarrow$ Suggerimento}\normalsize\color{black}
\begin{lstlisting}
(defun evaluate-expression (expr)
  ;; atomo numerico: restituisci direttamente
  (cond ((numberp expr) expr)
        ((symbolp expr) (error "variabile non supportata: ~A" expr))
        ;; lista: (operatore arg1 arg2)
        (t (let ((op  (car expr))
                 (a1  (evaluate-expression (cadr expr)))
                 (a2  (evaluate-expression (caddr expr))))
             (case op
               (+ (+ a1 a2))
               (- (- a1 a2))
               (* (* a1 a2))
               (/ (if (= a2 0)
                      (error "divisione per zero")
                      (/ a1 a2)))
               (otherwise (error "operatore sconosciuto: ~A" op)))))))
\end{lstlisting}
\end{solbox}

\label{sol:144}
\begin{solbox}{144}
\small\color{black!55}\hyperref[ex:144]{$\leftarrow$ Esercizio}\enspace\textcolor{black!30}{|}\enspace\hyperref[hint:144]{$\leftarrow$ Suggerimento}\normalsize\color{black}
\begin{ragionamentoBox}{144}
L'ambiente e' una alist \fn{((simbolo . valore) ...)}.
Il valutatore e' un'estensione dell'es.~143 con due nuovi casi:
simbolo (lookup) e forma \fn{set} (binding).
L'ambiente e' passato per valore; per modificarlo
permanentemente, \fn{set} deve restituire il nuovo ambiente
oppure usare un contenitore mutabile.
\end{ragionamentoBox}
\begin{lstlisting}
(defun mio-eval (expr env)
  "Valuta EXPR nell'ambiente ENV (alist simbolo->valore).
   Restituisce (valori . nuovo-env)."
  (cond
    ;; numero: restituisce se stesso
    ((numberp expr) (cons expr env))
    ;; simbolo: lookup nell'ambiente
    ((symbolp expr)
     (let ((binding (assoc expr env)))
       (if binding
           (cons (cdr binding) env)
           (error "variabile non definita: ~A" expr))))
    ;; forma (set var val)
    ((and (listp expr) (eq (car expr) 'set))
     (let* ((var (cadr expr))
            (res (mio-eval (caddr expr) env))
            (val (car res))
            (e2  (cdr res)))
       (cons val (acons var val e2))))
    ;; espressione aritmetica
    ((listp expr)
     (let* ((op (car expr))
            (r1 (mio-eval (cadr expr)  env))
            (r2 (mio-eval (caddr expr) (cdr r1)))
            (a  (car r1))
            (b  (car r2))
            (e2 (cdr r2)))
       (cons (case op
               (+ (+ a b)) (- (- a b))
               (* (* a b)) (/ (/ a b)))
             e2)))
    (t (error "espressione non riconosciuta: ~A" expr))))
\end{lstlisting}
\end{solbox}

% =============================================================
\section*{Progetti}
% =============================================================

\label{sol:145}
\begin{solbox}{145}
\small\color{black!55}\hyperref[ex:145]{$\leftarrow$ Esercizio}\normalsize\color{black}
\begin{lstlisting}
(defun gioca ()
  (let ((segreto (1+ (random 100)))
        (tentativi 0))
    (format t "Indovina il numero (1-100):~%")
    (loop
      (format t "> ")
      (let ((guess (read)))
        (incf tentativi)
        (cond ((< guess segreto) (format t "Troppo basso!~%"))
              ((> guess segreto) (format t "Troppo alto!~%"))
              (t (format t "Esatto! Tentativi: ~D~%" tentativi)
                 (return)))))))
\end{lstlisting}
\end{solbox}

\label{sol:146}
\begin{solbox}{146}
\small\color{black!55}\hyperref[ex:146]{$\leftarrow$ Esercizio}\enspace\textcolor{black!30}{|}\enspace\hyperref[hint:146]{$\leftarrow$ Suggerimento}\normalsize\color{black}
\begin{ragionamentoBox}{146}
Separazione netta tra Model (griglia), Logic (verifica vittoria,
mossa valida) e IO (input/output). La funzione \fn{vinto-p}
controlla le 8 combinazioni vincenti hard-coded
--- semplice e sufficiente per una griglia $3\times3$.
\end{ragionamentoBox}
\begin{lstlisting}
;;; MODEL
(defun crea-griglia () (make-array '(3 3) :initial-element nil))

;;; LOGIC
(defun mossa-valida-p (g r c)
  (null (aref g r c)))

(defun fai-mossa (g r c giocatore)
  (setf (aref g r c) giocatore))

(defun vinto-p (g p)
  (or ;; righe
      (loop for r below 3 thereis
        (loop for c below 3 always (eq (aref g r c) p)))
      ;; colonne
      (loop for c below 3 thereis
        (loop for r below 3 always (eq (aref g r c) p)))
      ;; diagonali
      (loop for i below 3 always (eq (aref g i i) p))
      (loop for i below 3 always (eq (aref g i (- 2 i)) p))))

(defun pareggio-p (g)
  (loop for r below 3 always
    (loop for c below 3 always (not (null (aref g r c))))))

;;; IO
(defun stampa-griglia (g)
  (dotimes (r 3)
    (dotimes (c 3)
      (format t "~A " (or (aref g r c) ".")))
    (terpri)))

;;; LOOP DI GIOCO
(defun gioca-tris ()
  (let ((g (crea-griglia))
        (turno 0))
    (loop
      (stampa-griglia g)
      (let ((p (if (evenp turno) 'X 'O)))
        (format t "Turno di ~A. Riga colonna: " p)
        (let ((r (read)) (c (read)))
          (if (mossa-valida-p g r c)
              (progn
                (fai-mossa g r c p)
                (cond ((vinto-p g p)
                       (stampa-griglia g)
                       (format t "~A vince!~%" p)
                       (return))
                      ((pareggio-p g)
                       (stampa-griglia g)
                       (format t "Pareggio!~%")
                       (return))
                      (t (incf turno))))
              (format t "Mossa non valida.~%")))))))
\end{lstlisting}
\end{solbox}

\label{sol:147}
\begin{solbox}{147}
\small\color{black!55}\hyperref[ex:147]{$\leftarrow$ Esercizio}\enspace\textcolor{black!30}{|}\enspace\hyperref[hint:147]{$\leftarrow$ Suggerimento}\normalsize\color{black}
\begin{lstlisting}
(defun conta-vicini (g r c rows cols)
  ;; conta le 8 celle adiacenti vive
  (let ((count 0))
    (loop for dr from -1 to 1 do
      (loop for dc from -1 to 1
            unless (and (= dr 0) (= dc 0))
            do (let ((nr (+ r dr)) (nc (+ c dc)))
                 (when (and (>= nr 0) (< nr rows)
                            (>= nc 0) (< nc cols)
                            (= (aref g nr nc) 1))
                   (incf count)))))
    count))

(defun evolvi (griglia)
  (let* ((rows (array-dimension griglia 0))
         (cols (array-dimension griglia 1))
         (nuova (make-array (list rows cols) :initial-element 0)))
    (dotimes (r rows nuova)
      (dotimes (c cols)
        (let ((v  (aref griglia r c))
              (n  (conta-vicini griglia r c rows cols)))
          (setf (aref nuova r c)
                (if (= v 1)
                    (if (member n '(2 3)) 1 0)
                    (if (= n 3) 1 0))))))))

(defun stampa-griglia-vita (g)
  (let ((rows (array-dimension g 0))
        (cols (array-dimension g 1)))
    (dotimes (r rows)
      (dotimes (c cols)
        (format t "~A" (if (= (aref g r c) 1) "#" ".")))
      (terpri))
    (terpri)))
\end{lstlisting}
\end{solbox}

\label{sol:148}
\begin{solbox}{148}
\small\color{black!55}\hyperref[ex:148]{$\leftarrow$ Esercizio}\enspace\textcolor{black!30}{|}\enspace\hyperref[hint:148]{$\leftarrow$ Suggerimento}\normalsize\color{black}
\textit{Soluzione di riferimento --- scheletro minimale per studenti.}

\begin{ragionamentoBox}{148}
Struttura a tre strati: MODEL (defstruct), LOGIC (operazioni),
IO (file e interazione). Il salvataggio usa \fn{print} (S-expression)
per semplicita': e' leggibile con \fn{read} senza parsing custom.
\end{ragionamentoBox}
\begin{lstlisting}
;;; MODEL
(defstruct studente nome cognome voto)

;;; ARCHIVIO
(defparameter *db* (make-hash-table :test #'equal))

;;; LOGIC
(defun inserisci (nome cognome voto)
  "Aggiunge o aggiorna uno studente."
  (let ((id (format nil "~A ~A" nome cognome)))
    (setf (gethash id *db*)
          (make-studente :nome nome :cognome cognome :voto voto))
    id))

(defun cerca (nome cognome)
  "Cerca uno studente; segnala errore se non esiste."
  (let ((id (format nil "~A ~A" nome cognome)))
    (or (gethash id *db*)
        (error "Studente non trovato: ~A" id))))

(defun elimina (nome cognome)
  "Rimuove uno studente."
  (remhash (format nil "~A ~A" nome cognome) *db*))

(defun filtra-per-voto (minimo)
  "Restituisce gli studenti con voto >= MINIMO."
  (loop for s being the hash-values of *db*
        when (>= (studente-voto s) minimo) collect s))

(defun ordina-per-voto ()
  "Restituisce la lista degli studenti ordinata per voto."
  (sort (loop for s being the hash-values of *db* collect s)
        #'> :key #'studente-voto))

;;; IO
(defun salva (file)
  "Salva il database su FILE come S-expressions."
  (with-open-file (s file :direction :output :if-exists :supersede)
    (maphash (lambda (k v)
               (print (list k (studente-nome v)
                              (studente-cognome v)
                              (studente-voto v))
                      s))
             *db*)))

(defun carica (file)
  "Carica il database da FILE."
  (clrhash *db*)
  (with-open-file (s file :direction :input :if-does-not-exist nil)
    (when s
      (loop for record = (read s nil nil)
            while record
            do (inserisci (second record)
                          (third record)
                          (fourth record))))))
\end{lstlisting}
\end{solbox}


